\chapter{Metodologia operacional}
\label{cap:metodologia}

Este capítulo descreve o desenho que o laboratório já montou: corpus, escore, memória, alternativas simples e protocolo de comparação. É o \emph{como}, não o veredito. Números de desempenho entram só no Capítulo~\ref{cap:andamentos}, e mesmo lá como andamento, não como prova de hipótese.

\section{Recorte e corpus}

O recorte privilegiado é o saneamento de capital aberto, com a Sabesp (\texttt{SBSP3}) como caso âncora. O corpus \emph{strict} restringe matérias à empresa, com data, fonte e texto alinháveis no tempo. Copasa e Sanepar permanecem no desenho como extensão, não como amostra já fechada neste manuscrito.

A coleta é parte do problema científico \cite{neuenschwander2014webmedia}: raspagem de portais, normalização de data, deduplicação e exclusão de textos que não tratam de fato da companhia. Sem esse filtro, o índice mede a agenda da imprensa, não a empresa.

\section{Do texto ao índice}

Cada notícia recebe um escore de sentimento de domínio em português, na linha de \citeonline{santos2023finbert}. Sejam $P(\mathrm{pos})$ e $P(\mathrm{neg})$ as probabilidades atribuídas pelo classificador. O escore contínuo por matéria é
\[
d = P(\mathrm{pos}) - P(\mathrm{neg}).
\]
O neutro permanece nas probabilidades; não o descartamos à força.

Os escores de um mesmo dia se combinam em um impacto diário. A série da empresa ganha memória por média móvel exponencial (EWMA), analogia conceitual ao decaimento de \citeonline{shapiro2022fed}, não réplica do índice do Federal Reserve. Em forma resumida,
\[
\mathrm{ITI}_t = \alpha_{\mathrm{eff}} \cdot \mathrm{ITI}_{t-1} + (1 - \alpha_{\mathrm{eff}}) \cdot \mathrm{impacto}_t,
\]
em que $\alpha$ controla memória --- valor alto retém mais passado; valor baixo reage mais ao dia --- e \textbf{não} é o alfa de excesso de retorno da literatura financeira. Dias sem notícia aplicam decaimento sobre o estado anterior. O predictor principal previsto para a comparação semanal é o estado líquido no último dia útil da semana.

\section{Alternativas simples (B0--B3)}

Qualquer índice mais elaborado precisa se justificar contra medidas cruas extraídas das \emph{mesmas} notícias \cite{loughran2011liability}. Definimos quatro baselines internos; nenhum usa preço de mercado.

\begin{table}[htbp]
\caption{Baselines internos derivados do mesmo corpus de notícias.}
\label{tab:baselines}
\begin{tabular}{|p{0.12\textwidth}|p{0.78\textwidth}|}
\hline
\textbf{Código} & \textbf{Conteúdo} \\
\hline
B0 & Contagem de notícias no período \\
\hline
B1 & Média do escore $d$ \\
\hline
B2 & Média de $d$ ponderada por confiança do classificador \\
\hline
B3 & Impacto diário \emph{sem} memória EWMA \\
\hline
\end{tabular}

\vspace{0.4em}
{\footnotesize Fonte: elaboração própria, a partir do protocolo operacional do laboratório.}

\end{table}

\section{Protocolo de comparação}

O preço da ação entra só como variável alvo. O alinhamento é \textbf{semanal}: o ITI da semana confronta a soma dos log-retornos diários nas $h$ semanas seguintes, com $h \in \{1,2,4\}$, na linha de horizontes múltiplos da evidência brasileira \cite{duarte2020news,eramiars2025correlacao}.

Há quatro baselines, duas métricas de associação (Pearson e Spearman) e três horizontes --- vinte e quatro confrontos cabeça a cabeça por execução. ``Vitória'' significa apenas que o ITI correlaciona \emph{melhor} que aquele baseline naquela métrica e naquele horizonte. Significância prevista: bootstrap em bloco, não causalidade \cite{utfpr2024bert}.

Antes de interpretar associação com o mercado, o desenho exige um \textbf{gate} de qualidade do classificador no corpus de saneamento: concordância com anotação humana (acurácia da ordem de 70\% ou $\kappa \geq 0{,}40$). O modelo de \citeonline{santos2023finbert} é ferramenta, não garantia nesse setor.

O Quadro~\ref{tab:protocolo} resume o recorte operacional. A Figura~\ref{fig:pipeline} situa os elos em sequência.

\begin{table}[htbp]
\caption{Protocolo operacional do caso âncora (Sabesp).}
\label{tab:protocolo}
\begin{tabular}{|p{0.28\textwidth}|p{0.64\textwidth}|}
\hline
\textbf{Item} & \textbf{Especificação} \\
\hline
Empresa & Sabesp (\texttt{SBSP3}) \\
\hline
Classificador & FinBERT-PT-BR; escore $d$ \\
\hline
Índice & EWMA; estado líquido no fim da semana \\
\hline
Baselines & B0--B3 (somente notícias) \\
\hline
Alvo & Retorno futuro semanal, $h \in \{1,2,4\}$ \\
\hline
Comparação & Pearson e Spearman; 24 duelos \\
\hline
Incerteza & Bootstrap em bloco \\
\hline
Gate de qualidade & Concordância humana (acurácia ou $\kappa$) \\
\hline
\end{tabular}

\vspace{0.4em}
{\footnotesize Fonte: elaboração própria.}

\end{table}

\begin{figure}[htbp]
\centering
\begin{tikzpicture}[
  font=\small,
  box/.style={draw, rectangle, rounded corners=1pt, align=center, inner sep=6pt, minimum height=1.15cm, text width=2.15cm},
  arr/.style={->, thick}
]
\node[box] (c) {Coleta\\corpus};
\node[box, right=0.35cm of c] (s) {Escore $d$\\FinBERT-PT};
\node[box, right=0.35cm of s] (e) {EWMA\\ITI};
\node[box, right=0.35cm of e] (b) {B0--B3\\mesmas\\notícias};
\node[box, right=0.35cm of b] (r) {Duelos\\semanais};
\draw[arr] (c) -- (s);
\draw[arr] (s) -- (e);
\draw[arr] (s) -- (b);
\draw[arr] (e) -- (r);
\draw[arr] (b) -- (r);
\end{tikzpicture}
\caption{Fluxo operacional: das notícias aos duelos do ITI com alternativas simples.}
\label{fig:pipeline}

\vspace{0.4em}
{\footnotesize Fonte: elaboração própria.}
\end{figure}

Nada neste capítulo afirma que o ITI acrescente informação ao mercado. Afirma só o compromisso de método: coletar com critério, classificar com modelo de domínio, agregar com memória explícita e só então perguntar se o índice diz mais do que contar matérias.
